\documentclass[../../main.tex]{subfiles}
\begin{document}

\begin{flushright}
\footnotesize\textit{【自動文字起こし・要確認】}
\end{flushright}

{\bf [解]}
$k\in\mathbb{Z}_{\ge1}$とする．

\textbf{(1)} $y=\cos x$は区間内で単調減少で$-1\le\cos x\le1$である．一方，
\begin{align*}
0\le\frac1{x^2+1}\le1\iff-1\le-1+\frac2{x^2+1}\le1\iff-1\le\frac{1-x^2}{1+x^2}\le1
\end{align*}
である．したがって，$f(x)=\cos x-\dfrac{1-x^2}{1+x^2}$とすると，
\begin{align*}
f(2k\pi)\ge0，\qquad f((2k+1)\pi)\le0
\end{align*}
となるから，$f$が連続であることも用いて，$f(\alpha_k)=0\ (2k\pi\le\alpha_k\le(2k+1)\pi)$なる$\alpha_k$がある（$C_1$，$C_2$の共有点の存在）．$C_1$の接線$\ell_k$は
\begin{align*}
\ell_k:y=-\sin\alpha_k(x-\alpha_k)+\cos\alpha_k=g(x)
\end{align*}
で表せる．$f(\alpha_k)=0$から
\begin{align*}
\cos\alpha_k=\frac{1-\alpha_k^2}{\alpha_k^2+1}，\qquad\sin\alpha_k=\sqrt{1-\frac{(1-\alpha_k^2)^2}{(\alpha_k^2+1)^2}}=\frac{2\alpha_k}{\alpha_k^2+1}\ (>0)
\end{align*}
であり，
\begin{align*}
g(0)=-\frac{2\alpha_k}{\alpha_k^2+1}(-\alpha_k)+\frac{1-\alpha_k^2}{\alpha_k^2+1}=\frac{\alpha_k^2+1}{\alpha_k^2+1}=1
\end{align*}
から，$\ell_k$は$(0,1)$を通る．$\blacksquare$

\textbf{(2)} $h(x)=g(x)-\dfrac{1-x^2}{1+x^2}$とおく．
\begin{align*}
h(x)&=-\frac{2\alpha_k}{\alpha_k^2+1}x+1-\frac{1-x^2}{1+x^2}\\
&=-\frac{2\alpha_k}{\alpha_k^2+1}x+2-\frac2{1+x^2}\\
&=\frac{-2x}{(\alpha_k^2+1)(x^2+1)}(x-\alpha_k)(\alpha_kx-1)
\end{align*}

$2k\pi\le x,\alpha_k\le(2k+1)\pi$（$k\ge1$）では$x>0$，$\alpha_kx-1>0$だから，$h(x)$の符号は$-(x-\alpha_k)$の符号と一致し，
\begin{align*}
x<\alpha_k\text{の時，}h(x)>0，\qquad x=\alpha_k\text{で，}h(x)=0，\qquad\alpha_k<x\text{の時，}h(x)<0
\end{align*}
となり，たしかに$2k\pi\le x\le(2k+1)\pi$では$h(x)=0$はただ1つの解を持つ．つまり，$C_1$と$C_2$の共有点はただ1つである．$\blacksquare$

\bigskip
\noindent{\bf [解2]}
\textbf{(2)} $x=t$における$C_1$の接線は$y=-\sin t(x-t)+\cos t$であり，これより，切片の関数$Y(t)=t\sin t+\cos t$，$Y'(t)=t\cos t$から下表を得る．

\begin{center}
\begin{tabular}{c|ccccc}
$t$ & $2k\pi$ & & $\left(2k+\frac12\right)\pi$ & & $(2k+1)\pi$ \\
\hline
$Y'$ & & $+$ & $0$ & $-$ & \\
\hline
$Y$ & $1$ & $\nearrow$ & & $\searrow$ & $-1$ \\
\end{tabular}
\end{center}

$\alpha_k\ne2k\pi$であることと，$2k\pi$から$(2k+1)\pi$で$Y(t)=1$なる$t$がただ1つ（$\alpha_k$自身）であることから，これと(1)をあわせて，題意は示された．$\blacksquare$

\newpage

\end{document}
